\chapter{Marco teórico y Antecedentes }
\hrule \bigskip \vspace*{1cm}
%------------------------------------------------------------------------
La forma como colocar un algoritmo es mediante el
\verb"\usepackage{algorithmic}" y \verb"{algorithm}" este imprime de
la siguiente forma:

\bigskip
\begin{algorithm}
\caption{Mapeamiento}\label{mapeadoEVA} processo\_ID(Identificación de flags)\\
Require: Lista de ${1\ldots N}$ que contenga los ID de las clases
correspondientes(provenientes del $FM$).
\begin{algorithmic} [1]
\STATE Generar \emph{lista} a partir de pares correspondientes según
$FM$ \WHILE {SchemaB contenga alguna clase}
\IF{valorASIG(\emph{claseB}) $\geq$ parametro \verb"VAL"}\STATE
\emph{claseB} $\Longleftarrow$ siguiente clase de SchemaB \STATE
\emph{lista} $\Longleftarrow$ agregar los términos de \emph{claseB}
y su correspondiente \emph{claseA} \STATE valor
(\emph{lista$(A_{i},B_{j})$})=\verb"POS"=$1$ \ENDIF \ENDWHILE
\end{algorithmic}
\end{algorithm}

NOTA: Este package no viene incluido por default en el \LaTeX ni con
esta plantilla, pero si es de mucha utilidad, esta disponible en
internet así como muchas otras. Si desean incluir un nuevo
\verb"\usepackage{Nombre_Package}", solo deben agregarla en el
archivo \verb"unsa.cls" en una linea y ya estará disponible.

\begin{table}[h]
  \centering
  \begin{tabular}{|c|c|c|c|c|}
  \hline
  % after \\: \hline or \cline{col1-col2} \cline{col3-col4} ...
  Me & Sem & Lug & Pos & Gen\\
  & Cas & & &\\
  \hline
  \hline
  a & bf & sd & as & hj \\
  a & bf & sdff & fg & ert \\
  a & bf & as & fg & klj \\
  \hline
\end{tabular}
  \caption{Como hacer una tabla}\label{tab:demo}
\end{table}


\begin{figure}
\begin{center}
\includegraphics[angle=45, width=5cm]{Graficos/escuela}
\caption{Logotipo de la EPCC}
\end{center}
\end{figure}
